\documentclass[11pt]{article}
\usepackage[utf-8]{inputenc}
\usepackage[margin=1in]{geometry}
\usepackage{xcolor}
\usepackage{listings}
\usepackage{hyperref}

\lstset{
    basicstyle=\ttfamily\small,
    breaklines=true,
    backgroundcolor=\color{gray!10},
    frame=single
}

\definecolor{orange}{RGB}{255, 140, 0}

\title{\textcolor{orange}{\textbf{Zaika Frontend - Interview Guide}}}
\author{10 Simple Questions with Answers}
\date{}

\begin{document}

\maketitle

\section*{Project Overview}

\textbf{Zaika} = Recipe sharing app. Users login, see recipes, add recipes, search.

\textbf{Tech Stack:} React 18 + TypeScript + React Router + Tailwind CSS + Axios

\textbf{Folder Structure:}
\begin{lstlisting}
src/
├── components/        # Button, Input, Card (reusable)
├── pages/             # Landing, Dashboard pages
├── layouts/           # DashboardLayout (sidebar)
├── context/           # AuthenticationContext (global state)
├── hooks/             # useAuth, useRecipe (custom logic)
├── config/            # axios setup
└── @types/            # TypeScript interfaces
\end{lstlisting}

---

\section*{Question 1: What is React?}

\textbf{Answer:} React is a JavaScript library for building UIs with \textbf{components}. Think of components as LEGO blocks - you build small pieces (Button, Card, Input) and combine them into bigger pages.

\textbf{Why?} 
\begin{itemize}
    \item Reusable - one Button component used everywhere
    \item Fast - only updates what changed (virtual DOM)
    \item Easy to maintain - changes in one place
\end{itemize}

---

\section*{Question 2: What is useState?}

\textbf{Answer:} useState is a React hook that remembers things. Like a sticky note for your component.

\textbf{Example from Landing page:}
\begin{lstlisting}
const [state, setState] = useState({ 
  email: "", 
  password: "" 
});

// User types in email field
<input onChange={(e) => 
  setState({...state, email: e.target.value})
} />

// Component re-renders with new email value
\end{lstlisting}

---

\section*{Question 3: What is useContext?}

\textbf{Answer:} useContext is like a public bulletin board. Instead of passing data through 10 components, put it on the board and anyone can grab it.

\textbf{In Zaika:}
\begin{lstlisting}
// Create bulletin board
const AuthenticationContext = createContext();

// Put user data on it
<AuthenticationContext.Provider 
  value={{ user, onLogin, onLogout }}>
  <App />
</AuthenticationContext.Provider>

// Any component grabs it
const { user, onLogout } = useContext(AuthenticationContext);
<p>User: {user}</p>
<button onClick={onLogout}>Logout</button>
\end{lstlisting}

---

\section*{Question 4: What is React Router?}

\textbf{Answer:} React Router is a GPS for your app. When user visits "/dashboard", show DashboardLayout. When user visits "/dashboard/recipe/123", show RecipeDetails.

\textbf{Zaika Routes:}
\begin{lstlisting}
/              → Landing (login page)
/dashboard     → Home (recipe list)
/dashboard/addrecipe → Add Recipe form
/dashboard/recipe/:id → Recipe details
\end{lstlisting}

The \texttt{:id} is a wildcard. \texttt{/recipe/123} captures id="123".

---

\section*{Question 5: What is Outlet?}

\textbf{Answer:} Outlet is a placeholder where child pages render.

\begin{lstlisting}
DashboardLayout:
┌──────────────────────────────────┐
│ Header (Zaika Logo)              │
├────────────┬────────────────────┤
│ Sidebar    │ <Outlet />         │  ← Child page renders here
│ Fixed      │ (Home, AddRecipe)  │     while sidebar stays
└────────────┴────────────────────┘
\end{lstlisting}

Without Outlet, you'd copy-paste sidebar in every page!

---

\section*{Question 6: Complete Authentication Flow}

\textbf{Step-by-step:}

\begin{enumerate}
    \item User opens app → Landing page checks: "Do you have token?"
    \item User enters email + password → clicks Login
    \item Frontend sends POST to \texttt{/auth/join}
    \item Backend validates → returns token
    \item Frontend saves to \texttt{sessionStorage}
    \item Updates AuthenticationContext with user email
    \item Redirects to /dashboard
    \item Dashboard shows because token exists
\end{enumerate}

\textbf{If token expires:} Axios interceptor catches 401 error → clears sessionStorage → redirects to login

---

\section*{Question 7: What is sessionStorage?}

\textbf{Answer:} Browser's temporary memory. Like a clipboard that:
\begin{itemize}
    \item Stores token after login
    \item Clears when you close the tab
    \item Used to check if user is logged in
\end{itemize}

\begin{lstlisting}
// After login
sessionStorage.setItem("token", data.token);
sessionStorage.setItem("email", data.email);

// On every request, axios adds:
Authorization: "Bearer " + sessionStorage.getItem("token");

// Logout
sessionStorage.clear();
\end{lstlisting}

---

\section*{Question 8: What are Axios Interceptors?}

\textbf{Answer:} Interceptors are guards that check every request and response.

\begin{lstlisting}
// BEFORE sending request
→ Add token to header automatically

// AFTER getting response
→ If success (200): Show "Success!" toast
→ If error (401): Clear sessionStorage, redirect to login

Result: No need to handle auth in every API call!
\end{lstlisting}

---

\section*{Question 9: TypeScript Basics}

\textbf{Answer:} TypeScript catches mistakes before code runs.

\begin{lstlisting}
// JavaScript (might fail at runtime)
const user = "john";
user.toUppercase();  // Wrong! But runs...

// TypeScript (catches before running)
const user: string = "john";
user.toUppercase();  
// ERROR! "toUpperCase" doesn't exist

// Define shape of data
type _STATE = {
  email: string;
  password: string;
};

const [state, setState] = useState<_STATE>({...});
setState({ email: 123 });  // ERROR! Should be string
\end{lstlisting}

---

\section*{Question 10: Tailwind CSS}

\textbf{Answer:} Tailwind is "CSS in class names". No separate CSS files.

\begin{lstlisting}
// Old way: Write separate CSS file
.login-button {
  background: orange;
  color: white;
  padding: 8px 24px;
}

// Tailwind way: All in className
<button className="bg-orange-500 text-white py-2 px-6 
                   hover:bg-orange-600">
  Login
</button>

// Responsive design
<div className="w-full md:w-1/2 lg:w-1/3">
  {/* 100% width on mobile, 50% on tablet, 33% on desktop */}
</div>
\end{lstlisting}

\textbf{Benefits:}
\begin{itemize}
    \item Faster - write HTML + styles together
    \item Consistent - predefined colors (orange-500, zinc-900)
    \item Smaller - only includes used styles
\end{itemize}

---

\section*{Interview Tips}

\begin{enumerate}
    \item \textbf{Q1-3:} React basics (components, hooks, state)
    \item \textbf{Q4-6:} Routing and navigation
    \item \textbf{Q7-8:} Authentication and API
    \item \textbf{Q9-10:} TypeScript and styling
\end{enumerate}

\textbf{Be ready to:}
\begin{itemize}
    \item Draw user login flow on whiteboard
    \item Explain what happens when token expires
    \item Show how useState works with form inputs
    \item Explain why Context avoids prop drilling
    \item Compare Tailwind vs traditional CSS
\end{itemize}

\end{document}
